% section: 付録の地図を閉じる

\section{付録の地図を閉じる}

この付録では、漸化式を次のような一つの流れとして眺めた。

\[
  \begin{array}{c}
    \text{隣り合う項の差}        \\
    \downarrow            \\
    \text{差分と離散的な積分}      \\
    \downarrow            \\
    \text{反復と力学系}         \\
    \downarrow            \\
    \text{微分方程式と数値計算}     \\
    \downarrow            \\
    \text{ニュートン法と近似}      \\
    \downarrow            \\
    \text{行列と固有値}         \\
    \downarrow            \\
    \text{母関数と組合せ}        \\
    \downarrow            \\
    \text{三角関数とチェビシェフ多項式} \\
    \downarrow            \\
    \text{不動点と漸近解析}       \\
    \downarrow            \\
    \text{確率・数論・アルゴリズム}
  \end{array}
\]

最初は、漸化式は数列の項を計算するための式に見える。
しかし、項と項の関係を丁寧に見ると、変化、反復、近似、対称性、線形性、極限という数学の基本概念が姿を現す。

漸化式を解くとは、一般項を一つ書くことだけではない。
その規則が何を記録し、どこへつながり、どの情報を残しているのかを読むことである。
式を見たときに、すぐ計算を始めるのではなく、
\[
  \text{この規則を、別の数学の言葉で言い換えられないか}
\]
と考えてみる。
そこから先に、漸化式の本当の広がりがある。
